% The smallest useful document with the quran package.
% Compile with: xelatex quran-getting-started.tex
\documentclass{article}
\usepackage[a5paper, margin=1.5cm]{geometry}
\usepackage{arabxetex}            % right-to-left Arabic via \begin{arab}[utf]
\newfontfamily\arabicfont[Script=Arabic, Scale=1.15]{Amiri}
\usepackage[uthmani]{quran}       % Uthmani script; default is the simple script

\parindent=0pt
\pagestyle{empty}

\begin{document}

\section*{Surah \surahname[1] (\textarab[utf]{\surahname*[1]})}
% One paragraph per ayah, Basmalah added automatically.
\begin{arab}[utf]
  \quransurah[1]
\end{arab}

\section*{Ayat al-Kursi (2:255)}
% A single ayah. The surah may be given by number or by name: [Al-Baqara].
\begin{arab}[utf]
  \quranayah[Al-Baqara][255]
\end{arab}

\section*{The last three surahs, as one paragraph}
% The starred form runs the ayahs together instead of one paragraph each.
\begin{arab}[utf]
  \quransurah*[112-114]
\end{arab}

\end{document}
